\documentclass[9pt]{article}
\usepackage[utf8]{inputenc}
\usepackage[margin=0.45in]{geometry}
\usepackage{enumitem}
\usepackage[hidelinks]{hyperref}
\usepackage{parskip}
\setlength{\parindent}{0pt}

\renewcommand{\familydefault}{\sfdefault}
\newcommand{\sect}[1]{\vspace{4pt}\noindent\textbf{\large #1}\vspace{3pt}\hrule\vspace{6pt}}

\begin{document}

% ---------- Header ----------
\begin{center}
  {\Large \textbf{Thirulok Sundar Mohan Rasu}}\\[3pt]
  Ann Arbor, MI 48105 $\,|\,$ +1 (248) 704-2911 $\,|\,$
  \href{mailto:thirulok@umich.edu}{thirulok@umich.edu} $\,|\,$
  \href{https://thiruloksundar.github.io}{thiruloksundar.github.io} $\,|\,$
  \href{https://github.com/Thiruloksundar}{github.com/Thiruloksundar} $\,|\,$
  \href{https://www.linkedin.com/in/thirulok-sundar-mohanrasu/}{linkedin.com/in/thirulok-sundar-mohanrasu}
\end{center}
\vspace{-3mm}

% ---------- Education ----------
\sect{Education}
\textbf{University of Michigan --- Ann Arbor} \hfill Aug 2025 -- Dec 2026\\
M.S., Electrical and Computer Engineering \hfill GPA: 4.0 / 4.0\\
Relevant courses: \emph{Algorithmic Robotics, Mobile Robotics, Matrix Methods for ML \& Signal Processing, Probability \& Random Processes}\\[5pt]
\textbf{Indian Institute of Technology (ISM) Dhanbad} \hfill Dec 2021 -- May 2025\\
B.Tech., Electrical Engineering. Member: \emph{CyberLabs (Machine Learning)} \hfill GPA: 8.30 / 10

% ---------- Skills ----------
\sect{Skills}
\textbf{Languages:} Python, SQL, C++, Julia, MATLAB\\
\textbf{ML \& Data:} Machine Learning, Feature Engineering, Pandas, NumPy, PySpark, Time-Series Forecasting, Data Ingestion \& Consolidation, JSON/NoSQL, AWS S3\\
\textbf{Robotics \& CV:} ROS2, Nav2, SLAM, LiDAR, State Estimation, Object Detection \& Segmentation, MuJoCo, Gazebo\\
\textbf{Tools:} PyTorch, OpenCV, SciPy, HuggingFace, GitHub, Linux

% ---------- Experience ----------
\sect{Experience}

\textbf{Electric Vehicle Center --- University of Michigan (Prof.\ Ziyou Song)} \hfill Oct 2025 -- Present\\
\emph{Graduate Research Assistant}
\vspace{-2mm}
\begin{itemize}[leftmargin=*,noitemsep,topsep=1pt]
  \item Built \textbf{data ingestion and processing pipelines} in \textbf{Python/Pandas} to collect, clean, and consolidate raw \textbf{electrochemical time-series datasets}; conducted \textbf{feature sensitivity (OAT) analysis} identifying key impedance parameters (SEI thickness, CPE exponent) from a 6-dimensional feature space across 117 EIS spectra
  \item Performed \textbf{feature engineering and PCA-based selection} on early-cycle battery data --- identified discharge capacity, Ah throughput, C-rates, DoD, and DVA features as key predictors; trained \textbf{attention-enhanced Seq2Seq} model achieving \textbf{2--5\% MAPE} on capacity-fade trajectories across 225 cells
\end{itemize}

\vspace{-1mm}
\textbf{HDR Lab --- University of Michigan} \hfill Jan 2026 -- Present\\
\emph{Graduate Research Assistant}
\vspace{-2mm}
\begin{itemize}[leftmargin=*,noitemsep,topsep=1pt]
  \item Developing real-time \textbf{3D depth estimation} algorithms for visuo-tactile sensors; built \textbf{iOS app} for automated data collection enabling large-scale dataset generation for model training on robotic arm
\end{itemize}

\vspace{-1mm}
\textbf{Carnegie Mellon University (Dr.\ Arun Balajee Vasudevan)} \hfill Jan 2024 -- Sep 2025\\
\emph{Research Intern}
\vspace{-2mm}
\begin{itemize}[leftmargin=*,noitemsep,topsep=1pt]
  \item Engineered \textbf{end-to-end data generation and labeling pipeline} to curate a \textbf{1M+ sample dataset} of real/synthetic audio pairs stored in structured \textbf{JSON (NoSQL) format}; fine-tuned \textbf{CLAP} audio-language model achieving \textbf{90\% accuracy} on deepfake speech detection; paper submitted to \textbf{Interspeech 2026} (under review) \emph{\textbf{\href{https://drive.google.com/file/d/1hdSPxQQA-PKayYQPPESZb11ofm7XrQRX/view?usp=sharing}{(preprint)}}}
\end{itemize}

\vspace{-1mm}
\textbf{Mowito} \hfill Mar 2023 -- Jul 2023\\
\emph{Robotics Software Intern}
\vspace{-2mm}
\begin{itemize}[leftmargin=*,noitemsep,topsep=1pt]
  \item Developed \textbf{Multi-head Mask R-CNN} models for \textbf{instance segmentation} and object detection in warehouse environments; integrated into production \textbf{robot perception pipelines} for autonomous navigation
  \item Processed and labeled robot vision data stored in \textbf{AWS S3 buckets}; built modular \textbf{data annotation pipeline} supporting multiple formats; implemented \textbf{perceptual hashing} for dataset deduplication (\textbf{30\% size reduction})
\end{itemize}

% ---------- Selected Projects ----------
\sect{Selected Projects}

\textbf{Semantic LiDAR-SLAM with Language-Guided Navigation} \href{https://github.com/chittinvarunav/VLSGN_ROB530-Team21}{(GitHub)} \hfill 2026\\
Integrated \textbf{Cartographer SLAM} with \textbf{Grounding DINO} in ROS2; DINO populates a persistent semantic map enabling \textbf{natural language commands} to trigger \textbf{Nav2} navigation goals --- zero-shot object-goal navigation without environment-specific training.

\textbf{VLA Model Fine-tuning for Robot Manipulation} \href{https://github.com/Thiruloksundar/OpenVLA-Libero-finetuning}{(GitHub)} \hfill 2025\\
\textbf{Fine-tuned OpenVLA} (7B vision-language-action model) on \textbf{LIBERO manipulation benchmarks} via \textbf{LoRA}; deployed on UMich HPC (NVIDIA A40). Training loss: 19.08 $\rightarrow$ 1.16 over 3 epochs.

\textbf{Robot State Estimation: EKF \& Particle Filters} \href{https://github.com/Thiruloksundar/robot_state_estimation}{(GitHub)} \hfill 2025\\
Implemented and benchmarked \textbf{Extended Kalman Filter} and \textbf{Particle Filter} for PR2 robot localization; analyzed kinematics model accuracy vs.\ compute trade-offs across five scenarios.

\textbf{ML Pipeline for Battery Aging Prediction} \hfill 2025\\
Built end-to-end \textbf{data processing and ML pipeline} in Python: ingested raw NMC cycling data, extracted features (discharge capacity, Ah throughput, C-rate, DoD) with \textbf{PCA-based selection}, and trained \textbf{attention-enhanced Seq2Seq} for multi-horizon capacity-fade forecasting; visualized degradation trajectories across 225 cells using Matplotlib/Seaborn.

\end{document}
